¿Cuál es el objetivo de la Ingeniería Inversa?

La ingeniería inversa busca desarmar y analizar un sistema para entender su funcionamiento interno y encontrar vulnerabilidades, con el fin de explotarlas en el caso de ciberdelincuentes, o bien para fortalecer las defensas del sistema y comprender el funcionamiento del malware en un equipo atacado~\cite{II}.

La ingeniería inversa se usa para analizar el malware para comprender cómo infecta, cómo persiste y cómo se comunica con el atacante; también se utiliza para detectar fallos de seguridad en el software (como errores de autenticación o vulnerabilidades en la lógica del sistema) y evaluar qué tan robusto es protegiendo credenciales o llaves de API. Además, fuera del área de seguridad, se emplea para comprender sistemas heredados y diseñar integraciones compatibles.


